\documentclass[12pt,twoside]{report}
\usepackage{amsmath}
\usepackage{amsfonts}
\usepackage{listings}

% some definitions for the title page
\newcommand{\reporttitle}{Real-time reconstruction of non-rigid objects from RGBD data}
\newcommand{\reportauthor}{Amelia Gordafarid Crowther}
\newcommand{\supervisor}{Name of supervisor}
\newcommand{\reporttype}{Type of Report/Thesis}
\newcommand{\degreetype}{Computer Science BEng} 

% load some definitions and default packages
%%%%%%%%%%%%%%%%%%%%%%%%%%%%%%%%%%%%%%%%%
% University Assignment Title Page 
% LaTeX Template
% Version 1.0 (27/12/12)
%
% This template has been downloaded from:
% http://www.LaTeXTemplates.com
%
% Original author:
% WikiBooks (http://en.wikibooks.org/wiki/LaTeX/Title_Creation)
%
% License:
% CC BY-NC-SA 3.0 (http://creativecommons.org/licenses/by-nc-sa/3.0/)
% 
%
%%%%%%%%%%%%%%%%%%%%%%%%%%%%%%%%%%%%%%%%%
%----------------------------------------------------------------------------------------
%	PACKAGES AND OTHER DOCUMENT CONFIGURATIONS
%----------------------------------------------------------------------------------------
\usepackage[a4paper,hmargin=2.8cm,vmargin=2.0cm,includeheadfoot]{geometry}
\usepackage{textpos}
\usepackage{natbib} % for bibliography
\usepackage{tabularx,longtable,multirow,subfigure,caption}%hangcaption
\usepackage{fncylab} %formatting of labels
\usepackage{fancyhdr} % page layout
\usepackage{url} % URLs
\usepackage[english]{babel}
\usepackage{amsmath}
\usepackage{graphicx}
\usepackage{dsfont}
\usepackage{epstopdf} % automatically replace .eps with .pdf in graphics
\usepackage{backref} % needed for citations
\usepackage{array}
\usepackage{latexsym}
\usepackage[pdftex,pagebackref,hypertexnames=false,colorlinks]{hyperref} % provide links in pdf

\hypersetup{pdftitle={},
  pdfsubject={}, 
  pdfauthor={},
  pdfkeywords={}, 
  pdfstartview=FitH,
  pdfpagemode={UseOutlines},% None, FullScreen, UseOutlines
  bookmarksnumbered=true, bookmarksopen=true, colorlinks,
    citecolor=black,%
    filecolor=black,%
    linkcolor=black,%
    urlcolor=black}

\usepackage[all]{hypcap}


%\usepackage{color}
%\usepackage[tight,ugly]{units}
%\usepackage{float}
%\usepackage{tcolorbox}
%\usepackage[colorinlistoftodos]{todonotes}
% \usepackage{ntheorem}
% \theoremstyle{break}
% \newtheorem{lemma}{Lemma}
% \newtheorem{theorem}{Theorem}
% \newtheorem{remark}{Remark}
% \newtheorem{definition}{Definition}
% \newtheorem{proof}{Proof}


%%% Default fonts
\renewcommand*{\rmdefault}{bch}
\renewcommand*{\ttdefault}{cmtt}



%%% Default settings (page layout)
\setlength{\parindent}{0em}  % indentation of paragraph

\setlength{\parindent}{0em}  % indentation of paragraph

\setlength{\headheight}{14.5pt}
\pagestyle{fancy}
\renewcommand{\chaptermark}[1]{\markboth{\chaptername\ \thechapter.\ #1}{}} 
%\fancyhead[RO]{\sffamily \textbf{\thepage}} %Page no.in the right on even pages
%\fancyhead[LE]{\sffamily \textbf{\thepage}} %Page no. in the left on odd pages

\fancyfoot[ER,OL]{\thepage}%Page no. in the left on
                                %odd pages and on right on even pages
\fancyfoot[OC,EC]{\sffamily }
\renewcommand{\headrulewidth}{0.1pt}
\renewcommand{\footrulewidth}{0.1pt}
\captionsetup{margin=10pt,font=small,labelfont=bf}


%--- chapter heading

\def\@makechapterhead#1{%
  \vspace*{10\p@}%
  {\parindent \z@ \raggedright \sffamily
    \interlinepenalty\@M
    \Huge\bfseries \thechapter \space\space #1\par\nobreak
    \vskip 30\p@
  }}

%--- chapter heading

\def\@makechapterhead#1{%
  \vspace*{10\p@}%
  {\parindent \z@ \raggedright \sffamily
        %{\Large \MakeUppercase{\@chapapp} \space \thechapter}
        %\\
        %\hrulefill
        %\par\nobreak
        %\vskip 10\p@
    \interlinepenalty\@M
    \Huge\bfseries \thechapter \space\space #1\par\nobreak
    \vskip 30\p@
  }}

%---chapter heading for \chapter*  
\def\@makeschapterhead#1{%
  \vspace*{10\p@}%
  {\parindent \z@ \raggedright
    \sffamily
    \interlinepenalty\@M
    \Huge \bfseries  #1\par\nobreak
    \vskip 30\p@
  }}	
\allowdisplaybreaks


% load title page
\begin{document}
% Last modification: 2015-08-17 (Marc Deisenroth)
\begin{titlepage}

\newcommand{\HRule}{\rule{\linewidth}{0.5mm}} % Defines a new command for the horizontal lines, change thickness here

%----------------------------------------------------------------------------------------
%	LOGO SECTION
%----------------------------------------------------------------------------------------

\includegraphics[width = 4cm]{./figures/imperial}\\[0.5cm] 

\center % Center everything on the page
 
%----------------------------------------------------------------------------------------
%	HEADING SECTIONS
%----------------------------------------------------------------------------------------

\textsc{\LARGE \reporttype}\\[1.5cm] 
\textsc{\Large Department of Computing}\\[0.5cm] 
\textsc{\large Imperial College of Science, Technology and Medicine}\\[0.5cm] 

%----------------------------------------------------------------------------------------
%	TITLE SECTION
%----------------------------------------------------------------------------------------

\HRule \\[0.4cm]
{ \huge \bfseries \reporttitle}\\ % Title of your document
\HRule \\[1.5cm]
 
%----------------------------------------------------------------------------------------
%	AUTHOR SECTION
%----------------------------------------------------------------------------------------

\begin{minipage}{0.4\textwidth}
\begin{flushleft} \large
\emph{Author:}\\
\reportauthor % Your name
\end{flushleft}
\end{minipage}
~
\begin{minipage}{0.4\textwidth}
\begin{flushright} \large
\emph{Supervisor:} \\
\supervisor % Supervisor's Name
\end{flushright}
\end{minipage}\\[4cm]




%----------------------------------------------------------------------------------------


%----------------------------------------------------------------------------------------
%	DATE SECTION
%----------------------------------------------------------------------------------------

{\large \today} % Date, change the \today to a set date if you want to be precise


\vfill % Fill the rest of the page with whitespace
Submitted in partial fulfillment of the requirements for the \degreetype~of Imperial College London

\end{titlepage}




% page numbering etc.
\pagenumbering{roman}
\clearpage{\pagestyle{empty}\cleardoublepage}
\setcounter{page}{1}
\pagestyle{fancy}

%%%%%%%%%%%%%%%%%%%%%%%%%%%%%%%%%%%%
\begin{abstract}
Your abstract.
\end{abstract}

\cleardoublepage
%%%%%%%%%%%%%%%%%%%%%%%%%%%%%%%%%%%%
\section*{Acknowledgments}
Comment this out if not needed.

\clearpage{\pagestyle{empty}\cleardoublepage}

%%%%%%%%%%%%%%%%%%%%%%%%%%%%%%%%%%%%
%--- table of contents
\fancyhead[RE,LO]{\sffamily {Table of Contents}}
\tableofcontents 


\clearpage{\pagestyle{empty}\cleardoublepage}
\pagenumbering{arabic}
\setcounter{page}{1}
\fancyhead[LE,RO]{\slshape \rightmark}
\fancyhead[LO,RE]{\slshape \leftmark}

%%%%%%%%%%%%%%%%%%%%%%%%%%%%%%%%%%%%
\chapter{Introduction}

\begin{figure}[tb]
\centering
\includegraphics[width = 0.4\hsize]{./figures/imperial}
\caption{Imperial College Logo. It's nice blue, and the font is quite stylish. But you can choose a different one if you don't like it.}
\label{fig:logo}
\end{figure}

Figure~\ref{fig:logo} is an example of a figure. 

\section{Motivation}

The reconstruction of 3-dimensional scenes using RGB-D data has been the subject of a lot of research recently. Due to the wide range of applications in fields such as Robotics and Augmented Reality, it has gained much interest.

However, most existing systems are designed only to reconstruct static geometry. This leads to hard limitations in the adaptability of the system to a changing environment. Worse still, highly deformable or dynamic objects such as a moving person cannot be reconstructed at all by a static system.

There has been a lot of interest in techniques for non-rigid reconstruction, but despite it being a hot research topic, few implementations of proposed systems are widely available.

The aim of this project is to provide such an implementation that can reconstruct a rigid object from a  non-rigid scene in real-time.


%%%%%%%%%%%%%%%%%%%%%%%%%%%%%%%%%%%%
\chapter{Background}

\section{Kinect Camera}

The first and most fundamental element of the system is the sensor used to obtain information about the surface to be reconstructed. The most commonly used sensor for this purpose is the Kinect camera. Originally designed as a commodity level camera for home games consoles, it was available directly to consumers. This means that Kinect cameras are widely and cheaply available to use.

However, they provide high quality information considering their price. This is provided to the system in the form of an ordinary RGB image, augmented with a fourth parameter indicating the perceived depth from the camera at that pixel. Using the available depth information, it is also possible to compute the normal to the surface at each pixel of the image. In some cases, depth information may be missing at certain pixels, producing so-called `holes' in the depth image. Errors such as these are not difficult to overcome but must be handled.


The sensor is capable of producing an RGB-D image to the system in 640x480 resolution at up to 30Hz. This is important, as it determines the amount of information that can be added to the system at each update. Additionally, it places a hard limit on the update rate of the system, which can never exceed the sensor rate.

\section{Rigid Surface Reconstruction}
\subsection{KinectFusion}

KinectFusion provides one such method for determining correspondences between point clouds. It makes one very important assumption: that the surface it is reconstructing is rigid, and will not be deformed while it is being observed by the system. While this is a very restricted approach to the problem, this algorithm laid the groundwork for later non-rigid systems such as DynamicFusion.

\subsubsection{Iterative Closest Point algorithm}
This is a widely used algorithm in this field which has been optimised and adapted for many purposes. Its purpose is simply to take two point clouds and return the transformation which most closely aligns the two clouds.
Psuedo-code for the ICP algorithm is shown below:

\begin{lstlisting}[language=Python]
def ICP(P, Q):
    transform = identity
    while not aligned(P, Q):
        correspondence = []
        for p in P:
            p_t = apply(transform, p)
            q = closest point to p_t in Q
            correspondence += (p, q)
        transform = least_squares(correspondence)
    return transform
\end{lstlisting}

\subsubsection{Principles of operation}

Initially, a pre-processing phase occurs, where depth and normal maps are constructed from the input data. The system then enters a continuously updating loop with three stages.

\begin{enumerate}
\item \textbf{Sensor pose estimation}: KinectFusion uses the ICP algorithm to align its prediction of the surface with the actual measurements of the surface it receives. The transform that is returned is rigid, which functions perfectly for this algorithm, but requires adaptation for a non-rigid case. 

\item \textbf{Update surface reconstruction}: Here, the system receives input from the sensor, which it uses to create a truncated signed distance function (TSDF) for that live frame. Together with the TSDFs from previous frames, a composite TSDF is created to represent the system's model of the scene. This composite TSDF is discretised into a 2-dimensional array and stored in global GPU memory.

\item \textbf{Surface prediction}: At this stage KinectFusion will produce its own depth map by ray-casting against its own model of the scene. This produces a prediction of what it expects the scene to look like.
\end{enumerate}


\section{Non-rigid Surface Reconstruction}
\subsection{DynamicFusion}

DynamicFusion was the first system to achieve dense, dynamic scene reconstruction in real-time. Both VolumeDeform and KillingFusion were heavily inspired by or based on DynamicFusion.

\subsubsection{Warp field}

The warp field is essentially a set of deformation nodes, each of which is denoted:

$$\mathcal{N}^t_{\textbf{warp}}=\{ \textbf{dg}_v, \textbf{dg}_w, \textbf{dg}_{se3} \}$$

Where the parameters refer to:

\begin{itemize}
\item $\textbf{dg}_v$: The position vector of the deformation node

\item $\textbf{dg}_w$: The radius of influence of the node, which determines the strength of the application of the transformation to the node's neighbours

\item $\textbf{dg}_{se3}$: The transformation of the deformation node, represented using dual quaternions
\end{itemize}


\subsubsection{Dual quaternions}

3


\subsubsection{Dual quaternion blending}

KD-Tree

$$\mathcal{W}(x_c) = SE3(\textbf{DQB}(x_c))$$

$$\textbf{DQB}(x_c) = \frac{\sum_{k \in N(x_c)} \textbf{w}_k(x_c)\hat{\textbf{q}}_{kc}}{\Vert \sum_{k \in N(x_c)} \textbf{w}_k(x_c)\hat{\textbf{q}}_{kc} \Vert}$$

\subsubsection{Warp function}

6

\subsubsection{Surface fusion}

7

\noindent\textbf{Principles of operation}


There are three key steps that occur in each iteration of the DynamicFusion algorithm, and they each take place whenever a new live frame is produced:

\begin{enumerate}
\item \textbf{Warp field estimation}: The warp field is estimated by choosing parameters such that the error function below is minimised:

$$E(\mathcal{W}_t, \mathcal{V}, D_t, \mathcal{E}) = \textbf{Data}(\mathcal{W}_t, \mathcal{V}_t, D_t) + \lambda\textbf{Reg}(\mathcal{W}_t, \mathcal{E})$$

This function consists of two terms:
\begin{itemize}
\item $\textbf{Data}(\mathcal{W}_t, \mathcal{V}_t, D_t) = \sum\limits_{u \in \Omega} \psi_{data}(\hat{\textbf{n}}_u^\top (\hat{\textbf{v}}_u - \textbf{vl}_u))$

This term represents the ICP cost between the model and the live frame. This step is similar to KinectFusion, although it uses a non-rigid version of ICP.

\item $\textbf{Reg}(\mathcal{W}_t, \mathcal{E}) = \sum\limits_{i=0}^n\sum\limits_{j \in \mathcal{E}(i)}\alpha_{ij}\psi_{reg}(\textbf{T}_{ic}\textbf{dg}^j_v - \textbf{T}_{jc}\textbf{dg}^j_v)$

This term which regularises the warp field. This ensures that motion is as smooth as possible, and that deformation is as rigid as possible.
\end{itemize}

Additionally, the $\lambda$ parameter can be adjusted to determine the relative importance of each term to the error function as a whole.

$E$ is minimised using Gauss-Newton non-linear optimisation.

\item \textbf{Surface fusion}: Using the estimated warp field, both surfaces can be warped back into the canonical model. This means they occupy the same space relative to each other, so can be directly compared.

\item \textbf{Updating the canonical model}: Now that the live frame has been warped into the canonical model, points can be added or removed from the canonical model as appropriate.
\end{enumerate}


\subsection{VolumeDeform}

VolumeDeform is different to the other examined techniques in that it doesn't rely solely on depth information. It also uses the RGB data probided by the sensor to extract additional information about the live frame.\\


\noindent\textbf{Principles of operation}

VolumeDeform operates similarly to the other methods outlined, in that it involves a loop that repeatedly acquires new knowledge and integrates it with an existing model.

\begin{enumerate}
\item \textbf{Mesh extraction}: The signed distance field representing the system's model of the scene is used to generate a deformed 3D mesh using Marching Cubes
\item \textbf{Correspondence detection}: The mesh is rendered to produce a depth map which can be used to generate correspondences.

Additionally, SIFT features are detected in the live frame and stored by the system. Since SIFT features are robust to changes in transformation, they can be recognised again in later frames. This provides extra information to inform the process of correspondence detection.

\item \textbf{Deformation optimisation}: The deformation field is optimised so that it is consistent with the observed values of colour and depth.
\end{enumerate}


\section{SDF-2-SDF}

SDF-2-SDF is a technique for static reconstruction.

precursor to KillingFusion, and designed in a very similar way

operates more implicitly - no usage of point clouds, so data of depth image is abstracted into a SDF.

For each depth image that is processed by the system, a signed distance field is created to represent the target object as viewed in that frame.

this sdf is warped into the canonical SDF using a gradient descent scheme based on a custom error function.

the canonical sdf represents the systems current best reconstruction of the target object.

once the current frames sdf has been succsefully warped into the canonical sdf, then it is integrated into the canonical sdf to add any new geometry to the reconstruction.

this process continues for each frame of data.

\subsection{Signed distance field}

A signed distance field (SDF) is a mapping $\phi : \mathbb{R}^3 \rightarrow  \mathbb{R}$ which yields the distance from the surface of an object at a particular point in space. 
Points inside the object correspond to negative values; points outside the object correspond to positive values and points on the surface are mapped to zero.

One of the important properties of a signed distance field is that the magnitude of the gradient of the SDF is equal to one at all points. This is called the level set property and is useful as a regularizer.

additionally, the gradient of the sdf is equivalent to the normal of the surface at the given point.


\subsection{Gradient descent scheme}



simplest:

$$ E_{geom} = \frac{1}{2} \sum\limits_{\textrm{all voxels}} (\phi_n(\Psi) - \phi_{global})^2  $$


additionally, the normals can be constrained to be approximately 1. Originally the dot product of the two relevant normals was used

$$ E_{norm} = \sum\limits_{\textrm{all voxels}}(1 - \nabla \phi_n(\Psi) \cdot \nabla\phi_{global}) $$

These two error functions can be combined proportionally using a scaling factor $\alpha_{norm}$. This factor can take on values in the range $[0..1]$, but is usually set to zero to improve performance.

$$E = E_{geom} + \alpha_{norm}E_{norm}$$

$$ E'_{geom} = (\phi_n(\Psi) - \phi_{global}) \nabla \phi_n(\Psi)$$

as long as the magnitude of the maximum update remains above a set threshold, this process of gradient descent will continue.

\subsection{Canonical SDF}


This canonical SDF is adjusted incrementally as new data is processed each frame. 
at the end of each update to $\Psi$, the accurately posed SDF must be integrated into the canonical SDF. This is done using a weighted average of all previous SDFs

\begin{align*}
\Phi_{n+1}(\textbf{V}) &= \frac{W_n(\textbf{V})\Phi_n(\textbf{V}) + \omega_{n+1}(\textbf{V})\phi_{n+1}(\textbf{V})}{W_{n}(\textbf{V}) + \omega_{n+1}(\textbf{V}}\\
W_{n+1}(\textbf{V}) &= W_n(\textbf{V}) + \omega_n(\textbf{V}) \\
\end{align*}

Here, the voxel weighting is defined as:

\[
    \omega(\textbf{V}) = \left\{\begin{array}{lr}
    1, & \text{for } \phi_{true}(\textbf{V}) > - \eta\\
    0, & \text{otherwise}\\
    \end{array}\right\}
\]

$\eta$ is a hyper-parameter that determines the expected thickness of the object 


\section{KillingFusion}

KillingFusion operates very similarly to SDF-2-SDF, but introduces two other terms to the error function to enable it to estimate a non-rigid transformation.

\begin{align*}
E_{non-rigid}(\Psi) &= E_{data}(\Psi) + \omega_kE_{Killing}(\Psi) + \omega_sE_{level-set}(\Psi)\\
E'_{data}(\Psi) &= (\phi_n(\Psi) - \phi_{global}) \nabla_{\phi_n}(\Psi)\\
E'_{Killing}(\Psi) &= 2H_{uvw}(vec(J^{\top}_{\Psi}) vec\big(J_{\Psi})\big)\begin{pmatrix}1\\\gamma \end{pmatrix}\\
E'_{level-set}(\Psi) &= \frac{|\nabla_{\phi_n}(\Psi)| - 1}{|\nabla_{\phi_n}(\Psi)|_\epsilon}H_{\phi_n}(\Psi)\nabla_{\phi_n}(\Psi)
\end{align*}

This error function imposes three requirements on a potential solution:
\begin{itemize}
\item \textbf{Data term}: Similar to DynamicFusion, contains a component for non-rigid transformation.
\item \textbf{Killing condition}: This term requires the flow field to be a Killing vector field. If it satisifies this condition, then it will produce isometric motion Since a perfect isometric motion would impose only rigid transformations, this term only needs to be minimised so that the motion is approximately rigid. Additionally, the requirement has been weakened to make it computationally more efficient.
\item \textbf{Level set property}: The minimisation of this term ensures that the gradient of the SDF function remains approximately 1
\end{itemize}

The additional parameters $\omega_k$ and $\omega_s$ are used to adjust the relative importance of the three terms when calculating $E_{non-rigid}$.

This error function is also much more easily parallelised, leading to greatly increased performance.\\



\subsection{Deformation field}

The deformation field is a mapping $\Psi : \mathbb{R}^3 \rightarrow \mathbb{R}^3$ that maps a point in space to its deformation relative to the canonical SDF.

The deformation field is only defined over the volume being reconstructed. Any points queried outside the volume return a deformation of zero.

In practical terms, the volume is discretized into voxels of a fixed length. At each voxel, a single displacement vector $u : \mathbb{R}^3$ is stored to describe the deformation field at that point.

The main task of KillingFusion is to estimate the value of this deformation field, using the error functions outlined.

\subsection{Killing vector field}

\chapter{Implementation details}


\section{Signed distance field}

The signed distance field is ultimately produced from the depth image provided by the Kinect camera. 
Instead of using the depth data to produce a point cloud, an implicit technique is used instead.
We define the SDF as follows:

$$\phi_{\text{true}}(\textbf{V}) = D(\pi(\textbf{V})) - \textbf{V}_Z$$

\[
    \phi(\textbf{V}) = \left\{\begin{array}{lr}
    -1,                    & \text{for } \phi_{true}(\textbf{V}) < - \delta\\
    \phi_{true}(\textbf{V}) / \delta, & \text{for } -\delta \leq \phi_{true}(\textbf{V}) \leq \delta \\
    1,                     & \text{for } \phi_{true}(\textbf{V}) > \delta\\
    \end{array}\right\}
\]

\subsection{Projection}

Here, $D(x)$ is the depth data at the point $x$ in the image provided by the Kinect camera.

$\pi(\textbf{V})$ is the function which projects the centre of the voxel \textbf{V} into a 2-dimensional point in the image plane. There are two main choices for projection:
\begin{itemize}
\item \textbf{Perspective}: This provides a more realistic model from the point-of-view of the camera filming the footage. 
However, the projection requires more expensive matrix multiplication. Additionally, an expensive vector magnitude operation is required to calculate depth from the camera.
Perspective projection also leaves large areas of the volume which cannot be reconstructed due to being projected outside of the volume.

\item \textbf{Orthographic}: This method simply discards the z-component of the vector to determine where it lies in the depth image. This is much more efficient, as it also allows distance from the camera to be efficiently calculated directly from the z-component.
\end{itemize} 

In the end, orthographic projection was used. This is because the signed distance function is evaluated an extremely large number of times (numbers?) each iteration, so must be very efficient.

\subsection{Role of $\delta$}

As part of the implementation of the SDF, the true distance from the surface is divided by $\delta$ and clamped to the range $[-1 \ldots 1]$.

This has the effect of limiting the estimation of the deformation field to areas near the surface of the surface boundary of the SDF. 

If a point lies too far outside the surface, then it is simply designated as `somewhere outside' the surface and treated like all other points of its kind. The same is true for points far inside the surface.

This ensures that unnecessary calculations are not carried out for voxels which are ultimately irrelevant to the geometry of the surface.

in keeping with killingfusion paper, a value of delta was used ten times the voxel length.


\section{Voxel grid}

the volume to be reconstructed is divided up into a regular grid of voxels. resolution 5-10 millimetres - mention sensor error of kinect camera

a gradient descent scheme is used to iteratively align the current SDF to the canonical sdf. this process is repeated until the magnitude of the maximum update is less than a certain threshold (0.1 millimetres)

$$\Psi_{n+1} = \Psi_n - \alpha E(\Psi)$$

\subsection{Canonical SDF}

In practice, the SDFs produced can be sampled at each voxel throughout the volume and a simple test can be performed to determine whether or not to increment $\Phi_n(\textbf{V})$ and $W_n(\textbf{V})$

the canonical sdf samples previous sdfs

at each voxel, two values are stored: omega and phi.
if $\phi_{true}(\textbf{V}) > - \eta$, then omega is incremented and the value of $\phi_{true}$ is added to the total of phi.

then, if the distance needs to be queried, $\frac{\phi}{\omega}$ is returned.

\subsection{Deformation field}

deformation field also split along voxel grid.

single vector stored at each voxel.

since each frame is initialised with the previous frames deformation, only one deformation field needs to be created in memory. this field is continually updated and shared by all SDFs created.

%%%%%%%%%%%%%%%%%%%%%%%%%%%%%%%%%%%%

\chapter{Contributions}


\section{CPU version}

The code for this part of the project can be found at:

\begin{center}
\verb|github.com/millie-crowther/fusion|
\end{center}

\subsection{Multi-threading}

The design of the KillingFusion system is very well suited to parallelism, since the update of each voxel only depends on the evaluation of the error function at that voxel. 
Of course, given the size of the voxel grid involved, it is best to use the massive parallelism of a GPU to maximise performance.
However, it is still possible to parallelise the system on the CPU and achieve a moderate increase in performance. 

make a diagram for running time with different sized thread pools

\section{GPU version}

The code for this part of the project can be found at:

\begin{center}
\verb|github.com/millie-crowther/gpu_fusion|
\end{center}

\subsection{C++}

\subsection{CUDA}

Four entry points to CUDA code from host:

\begin{enumerate}
\item \textbf{Initialise}: Initialises the system, sets warp field to zero. samples first SDF into canonical sdf. allocates memory.

\item \textbf{Estimate} $\Psi$: Most important kernel launch. updates deformation field.

\item \textbf{Get Canonical SDF}: copies the canonical sdf from device to host so that it can be viewed.

\item \textbf{Clean up}: cleans up the system. de-allocates memory

\end{enumerate}


%%%%%%%%%%%%%%%%%%%%%%%%%%%%%%%%%%%%
\chapter{Experimental Results}


%%%%%%%%%%%%%%%%%%%%%%%%%%%%%%%%%%%%
\chapter{Conclusion}


%% bibliography
\bibliographystyle{apa}


\end{document}
